\documentclass[11pt]{article}

\usepackage[margin=1in]{geometry}
\usepackage{lmodern}
\usepackage{microtype}
\usepackage{hyperref}
\usepackage{booktabs}
\usepackage{array}
\usepackage{longtable}
\usepackage{listings}
\usepackage{xcolor}
\usepackage{enumitem}
\usepackage{titlesec}

\hypersetup{
colorlinks=true,
linkcolor=black,
urlcolor=blue
}

\titleformat{\section}{\large\bfseries}{\thesection}{0.75em}{}

\lstset{
basicstyle=\ttfamily\small,
backgroundcolor=\color{gray!8},
frame=single,
breaklines=true,
columns=fullflexible,
literate=
  {<=}{{$\leq$}}1
  {>=}{{$\geq$}}1
}

\title{\textbf{Zendesk Ticket Request Quality Feedback Automation}\\
\vspace{4pt}
\large Open Implementation Guide}

\date{}
\author{}

\begin{document}

\maketitle

\section{Overview}

This document describes a Zendesk workflow that improves support ticket quality by automatically providing feedback to requesters when a ticket receives a low internal rating.

The system introduces:

\begin{itemize}
\item A conditional multi-select field allowing agents to record why a request was difficult to resolve
\item An automated requester email containing guidance for submitting better tickets
\item Internal audit notes documenting automated feedback events
\item Safeguards to prevent duplicate notifications
\end{itemize}

This workflow can help support teams reduce follow-up requests for missing information and improve overall ticket resolution efficiency.

\section{Current State}

Many support environments allow agents to rate the quality of a ticket request after resolving the issue.

However, without structured feedback mechanisms:

\begin{itemize}
\item Requesters may repeatedly submit incomplete tickets
\item Agents must spend time requesting missing information
\item Support teams lack reporting on common submission issues
\end{itemize}

The workflow described in this document addresses these issues through automation and structured feedback.

\section{Proposed System}

\subsection{Conditional Feedback Field}

A multi-select field titled:

\textbf{Bad Rating Reason}

is displayed only when the internal ticket rating is low.

\subsubsection{Field Behavior}

\begin{center}
\begin{tabular}{ll}
\toprule
Property & Configuration \\
\midrule
Field Type & Multi-select \\
Visibility & Visible only when rating $\leq$ 2 \\
Required & Required when visible \\
\bottomrule
\end{tabular}
\end{center}

\subsubsection{Field Values}

\begin{center}
\begin{tabular}{ll}
\toprule
Value & Tag \\
\midrule
Not enough information & not\_enough\_information \\
Poor contact method & poor\_contact\_method \\
Did not submit ticket & did\_not\_submit\_ticket \\
Walk-Up & walk-up \\
No response & no\_response \\
Not listed above & not\_listed\_above \\
\bottomrule
\end{tabular}
\end{center}

These values allow agents to document the most common ticket submission issues.

\section{Automation Logic}

A Zendesk trigger runs when a ticket is solved and rated poorly.

\subsection{Trigger Conditions}

\begin{lstlisting}
Ticket Status = Solved
User rating <= 2
Bad Rating Reason is present
Tag requester_feedback_sent NOT present
\end{lstlisting}

\subsection{Trigger Actions}

\begin{lstlisting}
Send Email to Requester
Add Internal Note
Add Tag requester_feedback_sent
\end{lstlisting}

The tag ensures the notification is only sent once per ticket.

\section{Dynamic Requester Email Template}

Paste the following template into the Zendesk trigger email body.

\begin{lstlisting}
Hello {{ticket.requester.first_name}},

We’re always happy to help and want to make sure your issues get resolved
as quickly as possible.

Based on your most recent request, here are a few things that would help us
resolve similar issues faster next time:

{% if ticket.tags contains "not_enough_information" %}
- Please include more detail about the issue including the application
  name, error messages, and screenshots when possible.
{% endif %}

{% if ticket.tags contains "poor_contact_method" %}
- Please include the best way to reach you such as a phone number or
  extension.
{% endif %}

{% if ticket.tags contains "no_response" %}
- If we ask follow-up questions, please reply promptly so we can continue
  troubleshooting your request.
{% endif %}

{% if ticket.tags contains "walk-up" %}
- Please submit a web-form ticket when possible rather than making a
  walk-up request so the issue can be properly tracked and routed.
{% endif %}

{% if ticket.tags contains "did_not_submit_ticket" %}
- Please submit a ticket through the web form so the request can be
  documented and routed correctly.
{% endif %}

{% if ticket.tags contains "not_listed_above" %}
- Providing clear details and contact information helps resolve requests
  faster.
{% endif %}

Thank you for helping us improve the support process.

- Support Team
\end{lstlisting}

\newpage

\section{Dynamic Internal Audit Note}

Paste the following template into the trigger's internal note action.

\begin{lstlisting}
Requester feedback automation executed.

User rating recorded: 1 or 2

Selected guidance sent:

{% if ticket.tags contains "not_enough_information" %}
- Not enough information
{% endif %}

{% if ticket.tags contains "poor_contact_method" %}
- Poor contact method
{% endif %}

{% if ticket.tags contains "no_response" %}
- No response to follow-up
{% endif %}

{% if ticket.tags contains "walk-up" %}
- Walk-Up request
{% endif %}

{% if ticket.tags contains "did_not_submit_ticket" %}
- Did not submit ticket
{% endif %}

Tag applied: requester_feedback_sent
\end{lstlisting}

\section{Benefits}

\begin{itemize}
\item Improves ticket submission quality
\item Reduces resolution delays caused by missing information
\item Creates measurable data on ticket submission issues
\item Encourages structured communication between requesters and support teams
\end{itemize}

\section{Implementation Steps}

\subsection{Step 1: Create the Multi-Select Field}

\begin{enumerate}
\item Open Zendesk Admin Center
\item Navigate to \textbf{Objects and Rules → Tickets → Fields}
\item Click \textbf{Add Field}
\item Select \textbf{Multiselect}
\item Name the field \textbf{Bad Rating Reason}
\item Add the values listed in Section 3
\item Save the field
\end{enumerate}

\subsection{Step 2: Add the Field to the Ticket Form}

\begin{enumerate}
\item Navigate to \textbf{Objects and Rules → Tickets → Forms}
\item Open the active ticket form
\item Add the field \textbf{Bad Rating Reason}
\item Create a condition:
\end{enumerate}

\begin{lstlisting}
IF User rating <= 2
SHOW Bad Rating Reason
AND make required
\end{lstlisting}

\subsection{Step 3: Create the Automation Trigger}

\begin{enumerate}
\item Go to \textbf{Objects and Rules → Business Rules → Triggers}
\item Click \textbf{Add Trigger}
\item Name it \textbf{Requester Feedback Automation}
\item Configure conditions listed in Section 4
\item Configure actions listed in Section 4
\end{enumerate}

\subsection{Step 4: Paste Email Template}

Copy the template from:

\textbf{Section 5 – Dynamic Requester Email Template}

Paste it into the trigger's email action.

\subsection{Step 5: Paste Internal Note}

Copy the template from:

\textbf{Section 6 – Dynamic Internal Audit Note}

Paste it into the internal note action.

\subsection{Step 6: Test the Workflow}

\begin{enumerate}
\item Create a test ticket
\item Set rating = 1
\item Confirm the Bad Rating Reason field appears
\item Select a reason
\item Solve the ticket
\item Verify:
\begin{itemize}
\item Requester email was sent
\item Internal note was created
\item Automation tag was applied
\end{itemize}
\end{enumerate}

\section{Conclusion}

This automation provides a structured, scalable approach to improving support ticket quality while maintaining a positive and constructive tone with requesters.

By integrating conditional fields, dynamic messaging, and automated documentation, support teams can reduce resolution delays and improve operational efficiency.

\end{document}
